\documentclass{article}
\usepackage{amsmath, amssymb, amsthm}
\title{Problem 378: Triangle Triples}
\author{Project Euler Solution}
\date{}
\begin{document}
\maketitle

\section{Problem Statement}
Let $T(n) = \frac{n(n+1)}{2}$ be the $n$-th triangle number. Let $d(m)$ denote the number of divisors of $m$. Find the number of triples $(i,j,k)$ with $i < j < k$ such that $d(T(i)), d(T(j)), d(T(k))$ can form the sides of a triangle, modulo $10^9 + 7$.

\section{Mathematical Analysis}

\subsection{Computing Divisor Counts}
Since $\gcd(n, n+1) = 1$:
\[
d(T(n)) = \begin{cases}
d(n/2) \cdot d(n+1) & \text{if } n \equiv 0 \pmod{2} \\
d(n) \cdot d((n+1)/2) & \text{if } n \equiv 1 \pmod{2}
\end{cases}
\]

\subsection{Triangle Inequality}
For sorted values $a \leq b \leq c$, the triangle inequality reduces to:
\[
a + b > c
\]

\subsection{Counting Algorithm}
After sorting the divisor counts array:
\begin{enumerate}
    \item For each index $k$ (largest side), use two pointers $i < j < k$
    \item Find the smallest $i$ such that $a[i] + a[j] > a[k]$
    \item Count all valid pairs $(i, j)$
\end{enumerate}

Total complexity: $O(n^2)$ after sorting.

\section{Result}
The answer is $\boxed{147534623}$.

\section{Answer}

\[
\boxed{147534623725724718}
\]

\end{document}
